\section{Các công việc liên quan}\label{sec:related}

Mục này tổ chức theo chủ đề: tấn công tiêm lệnh, jailbreak, benchmark, phòng thủ (ba nhóm), các bài tổng quan, và cuối cùng là khoảng trống mà đề tài giải quyết một phần. Mô tả mỗi bài chỉ dựa trên những gì nhóm đã đối chiếu với trang bài báo hoặc bản PDF; mọi số liệu được nêu đều thuộc bài gốc, không phải của nhóm.

\subsection{Tấn công tiêm lệnh và khung hình thức}
\citet{perez2022ignore} là một trong những công trình sớm phân tích hai mục tiêu của tấn công bằng prompt: chiếm quyền mục tiêu (goal hijacking) và rò rỉ prompt (prompt leaking). Hai mục tiêu này tương ứng với các mục tiêu \texttt{integrity} và \texttt{leak} trong bộ prompt của nhóm. \citet{greshake2023indirect} chuyển trọng tâm sang tiêm lệnh \emph{gián tiếp}: lệnh nằm trong dữ liệu mà ứng dụng tự truy xuất. Bài đề xuất bốn cách đưa lệnh vào (thụ động, chủ động, do người dùng thực hiện và che giấu) cùng sáu nhóm hậu quả; đây là cơ sở cho mô hình đe doạ của nhóm (\cref{sec:threat}).

\citet{liu2024formalizing} đưa ra định nghĩa hình thức mà báo cáo này dùng làm khung: dữ liệu bị nhiễm $\tilde{x}$ được tạo từ dữ liệu đích $x_t$, chỉ dẫn tiêm vào $s_e$ và dữ liệu tiêm vào $x_e$. Bài phân tích năm kiểu tấn công cơ bản (Naive, Escape Characters, Context Ignoring, Fake Completion và Combined), mười cách phòng thủ chia thành nhóm phòng ngừa (prevention) và nhóm phát hiện (detection), cùng các thước đo ASV, MR, PNA-T, FPR và FNR. Cuộc thi HackAPrompt~\citep{schulhoff2023hackaprompt} thu thập một lượng lớn prompt tấn công từ người chơi và đề xuất một sơ đồ phân loại kỹ thuật chi tiết (ví dụ Context Ignoring, Few Shot, Refusal Suppression, Obfuscation, Payload Splitting, Defined Dictionary); nhóm dùng sơ đồ này để gắn nhãn kỹ thuật cho từng prompt (\cref{app:prompts}). Tensor Trust~\citep{toyer2024tensortrust} là bộ dữ liệu hơn 126.000 tấn công và 46.000 phòng thủ do người chơi một trò chơi trực tuyến tạo ra, kèm benchmark cho hai loại tấn công là trích xuất prompt và chiếm quyền.

\subsection{Jailbreak và các kỹ thuật nâng cao}
Tổng quan của \citet{xu2025survey}, dựa trên khung của \citet{liu2024formalizing}, phân biệt jailbreak với tiêm lệnh. Nhóm hiểu sự phân biệt này như sau: jailbreak nhắm vào huấn luyện an toàn của mô hình, còn tiêm lệnh nhắm vào việc mô hình làm theo chỉ dẫn đến từ một nguồn không có thẩm quyền; hai loại giao nhau ở kỹ thuật viết prompt. Nhóm vay mượn bốn họ kỹ thuật jailbreak để làm phong phú bộ prompt. Họ roleplay/DAN được \citet{shen2024dan} khảo sát trên các prompt thu thập ngoài thực tế (1.405 prompt jailbreak trong 15.140 prompt). \citet{wei2023jailbroken} giải thích vì sao jailbreak thành công qua hai kiểu thất bại của huấn luyện an toàn: mục tiêu cạnh tranh (competing objectives) và tổng quát hoá không khớp (mismatched generalization); nhóm dùng khái niệm thứ hai để lý giải nhóm tấn công che giấu (mã hoá, đổi hình thức chữ). \citet{yong2023lowresource} cho thấy việc dịch yêu cầu sang ngôn ngữ ít tài nguyên có thể né các cơ chế an toàn của GPT-4, là cơ sở cho nhóm tấn công đổi ngôn ngữ (code-switching). Cuối cùng, many-shot jailbreaking~\citep{anil2024manyshot} khai thác ngữ cảnh dài bằng cách nhồi nhiều ví dụ giả trước yêu cầu thật.

\subsection{Benchmark và bộ dữ liệu}
BIPIA~\citep{yi2025bipia} được giới thiệu là benchmark đầu tiên cho tiêm lệnh gián tiếp, kèm hai nhóm phòng thủ: hộp đen (không sửa mô hình, ví dụ nhắc nhở tường minh và đánh dấu ranh giới) và hộp trắng (có huấn luyện). InjecAgent~\citep{zhan2024injecagent} gồm 1.054 ca kiểm thử với 17 công cụ của người dùng và 62 công cụ của kẻ tấn công, tập trung vào tác tử dùng công cụ; trên ReAct-GPT-4 bài ghi nhận tấn công thành công 24\%. AgentDojo~\citep{debenedetti2024agentdojo} là môi trường động gồm 97 tác vụ và 629 ca kiểm thử bảo mật (theo bản nhóm đã đối chiếu; các phiên bản sau của bài có thể ghi số khác), đo ASR cùng với utility trên cùng một môi trường; đây là hình mẫu cho cách đo cặp (an ninh, tiện ích) mà nhóm theo. \Cref{tab:related-bench} so sánh các nguồn này với thiết lập của đề tài.

\begin{table}[t]
\centering\small
\caption[So sánh với các benchmark tiêu biểu]{So sánh thiết lập của đề tài với các benchmark tiêu biểu. Số liệu của các bài là theo bài gốc.}\label{tab:related-bench}
\begin{tabularx}{\linewidth}{@{}p{2.6cm}YYY@{}}
\toprule
 & Đối tượng & Tác vụ / quy mô & Điểm khác với đề tài\\
\midrule
BIPIA~\citep{yi2025bipia} & LLM đọc nội dung ngoài & Benchmark gián tiếp; phòng thủ hộp đen và hộp trắng & Tiếng Anh; không đo cộng dồn nhiều lớp\\
InjecAgent~\citep{zhan2024injecagent} & Tác tử dùng công cụ & 1.054 ca; 17 công cụ người dùng, 62 công cụ kẻ tấn công & Có công cụ đầy đủ\\
AgentDojo~\citep{debenedetti2024agentdojo} & Tác tử, 4 môi trường & 97 tác vụ, 629 ca bảo mật & Đo ASR cùng utility; môi trường lớn\\
Tensor Trust~\citep{toyer2024tensortrust} & Trò chơi chiếm quyền / trích xuất & $>$126 nghìn tấn công, 46 nghìn phòng thủ & Dữ liệu từ người chơi; không có tác vụ tóm tắt\\
\textbf{Đề tài} & Một LLM, tóm tắt/hỏi đáp tài liệu & Bộ prompt nhỏ tiếng Việt (\cref{sec:eval-data}) & Đo cộng dồn L1--L4, từ chối oan, độ trễ, hold-out\\
\bottomrule
\end{tabularx}
\end{table}

\subsection{Phòng thủ nhóm 1: lọc và phát hiện đầu vào}
Cách tiếp cận trực tiếp nhất là kiểm tra đầu vào trước khi đưa cho mô hình. \citet{alon2023perplexity} nhận thấy hậu tố đối kháng có perplexity rất cao khi chấm bằng GPT-2; lọc bằng perplexity đơn thuần cho nhiều dương tính giả trên prompt thường, và bài khắc phục bằng bộ phân loại dùng thêm độ dài token. \citet{jain2023baseline} khảo sát các phòng thủ nền tảng gồm phát hiện bằng perplexity, diễn đạt lại (paraphrase), tái phân đoạn token (retokenization) và huấn luyện đối kháng, cân nhắc giữa độ vững và cái giá về hiệu năng. Với tấn công viết bằng ngôn ngữ tự nhiên, như phần lớn prompt của đề tài, perplexity ít có tác dụng, nên L1 của nhóm dựa trên chuẩn hoá và quy tắc.

Nhiều bộ phát hiện là mô hình phân loại. Llama Prompt Guard~2~\citep{promptguard2} là bộ phân loại nhỏ (22M và 86M tham số) gán nhãn lành tính/độc hại cho prompt, nhưng yêu cầu xin quyền truy cập. GenTel-Safe~\citep{li2024gentel} gồm bộ phát hiện GenTel-Shield và benchmark GenTel-Bench (84.812 tấn công, 28 kịch bản). InjecGuard~\citep{li2024injecguard} chỉ ra hiện tượng \emph{over-defense}: các bộ chặn hiện hành nhạy với từ khoá thường gặp trong tấn công, nên độ chính xác trên mẫu lành tính chứa các từ đó giảm xuống khoảng 60\% (bộ NotInject có 339 mẫu). Đây là lý do tập benign của nhóm cố ý chứa các cụm từ như ``system prompt''. Attention Tracker~\citep{hung2024attentiontracker} phát hiện không cần huấn luyện bằng cách theo dõi sự dịch chuyển chú ý của một số đầu attention, nhưng cần truy cập bên trong mô hình nên không dùng được qua API thông thường. Cuối cùng, \citet{choudhary2025hownot} chỉ ra một lỗ hổng cấu trúc của phát hiện bằng câu trả lời đã biết (known-answer detection): kiểu tấn công thích ứng DataFlip đưa tỉ lệ phát hiện về 0\%. Đây là cảnh báo cần có khi dùng một LLM làm bộ chấm cho L2, vì bộ chấm cũng đọc đúng văn bản không tin cậy.

\subsection{Phòng thủ nhóm 2: tách chỉ dẫn khỏi dữ liệu}
Spotlighting~\citep{hines2024spotlighting} là họ kỹ thuật chỉnh prompt để mô hình phân biệt được các nguồn đầu vào, gồm đánh dấu ranh giới (delimiting), chèn dấu vào văn bản (datamarking) và mã hoá (encoding); với các mô hình GPT, bài báo cáo tỉ lệ tấn công thành công giảm từ trên 50\% xuống dưới 2\% với tác động tối thiểu đến hiệu năng tác vụ. L3 của nhóm dựa trên ý tưởng này, kết hợp thêm các phòng thủ prompt của \citet{liu2024formalizing} (sandwich, nhắc lại chỉ dẫn) và \citet{yi2025bipia} (nhắc nhở tường minh). Các bài sau đây thuộc hướng \emph{huấn luyện} và nhóm \textbf{không} tái hiện: Instruction Hierarchy~\citep{wallace2024hierarchy} huấn luyện GPT-3.5 để ưu tiên chỉ dẫn có đặc quyền cao hơn; StruQ~\citep{chen2025struq} tách prompt và dữ liệu thành hai kênh rồi tinh chỉnh mô hình bỏ qua chỉ dẫn nằm trong kênh dữ liệu; SecAlign~\citep{chen2025secalign} dùng tối ưu hoá ưa thích, báo cáo đưa tỉ lệ thành công của nhiều kiểu tiêm lệnh xuống dưới 10\% mà giữ tiện ích gần như cũ. CaMeL~\citep{debenedetti2025camel} thay đổi kiến trúc hệ thống để dữ liệu không chiếm được luồng điều khiển; trên AgentDojo bài báo cáo giải được 77\% tác vụ với bảo đảm an ninh có thể chứng minh, so với 84\% của hệ thống không phòng thủ.

\subsection{Phòng thủ nhóm 3: kiểm tra nhất quán và độ lệch tác vụ}
MELON~\citep{zhu2025melon} phát hiện tấn công bằng cách chạy lại quỹ đạo của tác tử với prompt người dùng \emph{bị che} trong khi giữ nguyên kết quả công cụ; nếu hai lần chạy cho hành động tương tự nhau thì hành vi không còn phụ thuộc vào tác vụ người dùng, tức dữ liệu đang điều khiển mô hình. Bài giới thiệu ba thiết kế để giảm dương tính giả và âm tính giả, và báo cáo trên AgentDojo vượt các phòng thủ so sánh cả về ngăn tấn công lẫn giữ tiện ích; đây là nguồn của L4. \citet{abdelnabi2025drift} tiếp cận từ bên trong mô hình: so sánh activation trước và sau khi đọc dữ liệu ngoài để phát hiện ``độ lệch tác vụ''; cách này cần truy cập activation nên nhóm chỉ nêu mà không cài đặt.

\subsection{Các bài tổng quan}
\citet{wang2026landscape} là bài hệ thống hoá (SoK) phân loại tấn công theo cách tạo payload (heuristic hoặc tối ưu hoá) và phòng thủ theo giai đoạn can thiệp (văn bản, mô hình, thực thi). Bài lưu ý nhiều phòng thủ và benchmark bỏ qua tác vụ phụ thuộc ngữ cảnh, và nhiều phòng thủ trông hiệu quả chỉ vì chúng triệt tiêu đầu vào ngữ cảnh. Quan sát này dẫn thẳng đến việc đo từ chối oan và utility cùng với ASR trong đề tài. \citet{xu2025survey} cung cấp bức tranh nền về tấn công LLM theo giai đoạn huấn luyện và suy luận.

\subsection{Khoảng trống và vị trí của đề tài}
Trong số tài liệu nhóm đã khảo sát, các công trình thường (i) đánh giá bằng tiếng Anh, (ii) tập trung vào tác tử có công cụ hoặc vào một phòng thủ riêng lẻ, hoặc (iii) không đo đồng thời hiệu quả, tỉ lệ từ chối oan, chất lượng và độ trễ trên cùng một bộ dữ liệu. Đề tài giải quyết \emph{một phần} khoảng trống đó ở quy mô nhỏ: đo cộng dồn bốn lớp trên một LLM chạy cục bộ với tác vụ tóm tắt tiếng Việt, kèm tập benign có chủ đích gây chặn nhầm và một tập hold-out để ước lượng khả năng tổng quát hoá. Nhóm không tuyên bố vượt các kết quả đã công bố trên mô hình lớn hơn.
